\subsection*{stacked triangles}

\label{sec:stacked_triangles}

\begin{center} \includegraphics [scale=0.5] {angle_bisector_r8.png} \end{center}

We stack two right triangles so that the hypotenuse of the first is the base of the second.  We draw the altitude to form another right angle.  

We immediately see an angle equality as marked by the black dots, most simply obtained by using vertical angles to see that we have similar right triangles.  Alternatively, compute sums, adding up the angles in the upper triangle compared to the one whose side is the altitude.

We draw a box and work our way through the diagram marking the other angles.

\begin{center} \includegraphics [scale=0.5] {angle_bisector_r9.png} \end{center}
 
 The relationships can be proven using the alternate interior angles theorem.

\begin{center} \includegraphics [scale=0.5] {angle_bisector_r10.png} \end{center}
